\documentclass[12pt,oneside]{article}
\usepackage{etoolbox}
\usepackage[T1]{fontenc}
\usepackage[full]{textcomp}
\usepackage{fancyhdr}
\usepackage{longtable}
%\usepackage[dvips]{graphics,color}      % usual driver for dvi or ps viewer
\usepackage[]{graphicx,color}      % usual driver for converting to pdf later
\usepackage[bottom=1in]{geometry}
\usepackage{verbatim}
%\usepackage[cmintegrals,cmbraces]{newtxmath}
\usepackage[intlimits]{amsmath}    % need for subequations
%\usepackage[]{amsmath}    % need for subequations
\usepackage{amsfonts}
\usepackage{amssymb}
%\usepackage{ebgaramond-maths}
\usepackage[sb]{ebgaramond} %semi-bold option
%\input{../ebgaramond-missing.tex}
\usepackage[boxed]{algorithm2e}
\usepackage{tikz}
\usepackage{float}
\usepackage{courier}
\usepackage{listings}
\usepackage{enumitem}
\usepackage{setspace}
\usepackage{caption}
\usepackage{longtable}
\usepackage{cellspace}
%\usepackage{multirow}
%\usepackage{subsection}

%\renewcommand{\rmdefault}{ugm} %old Garamond
\newcommand{\ttt}[1]{\texttt{\footnotesize{#1}}}
\graphicspath{./}
\DeclareGraphicsExtensions{.pdf,.eps}
\newcommand{\beqn}{\begin{equation}}
\newcommand{\eeqn}{\end{equation}}
\newcommand{\beqnn}{\begin{equation*}}
\newcommand{\eeqnn}{\end{equation*}}
\newcommand{\half}{\text{\scriptsize\ensuremath{\frac{1}{2}}}}
\newcommand{\redsize}[1]{\text{\scriptsize\ensuremath{#1}}}
\newcommand{\bsym}[1]{\boldsymbol{#1}}
\newcommand{\tinysize}[1]{\text{\tiny\ensuremath{#1}}}
\newcommand{\eee}{\text{\large{\rm e}}}
\newcommand{\E}{{\rm I\kern-.3em E}}
\newcommand{\pr}{\mathbb{P}}
\newcommand{\var}{{\rm Var}}
\newcommand{\defas}{\overset{*}{=}}
\newcommand{\sL}{{L}}
\newcommand{\edge}[2]{#1\!\!\rightarrow\!\!#2}
\newcommand{\sedge}[2]{#1\rightarrow#2}
\newcommand{\sumno}{\sum\nolimits}
\newcommand{\ns}[1]{\mbox{\normalsize $#1$}}

\let\osetminus\setminus
\renewcommand{\setminus}{\!\osetminus\!}
\let\oin\in

\renewcommand{\in}{\!\oin\!}

\numberwithin{equation}{section}
\renewcommand{\thefootnote}{\fnsymbol{footnote}}
\DeclareMathOperator*{\argmin}{arg\,min}
\DeclareMathOperator*{\argmax}{arg\,max}
\DeclareMathOperator*{\mino}{min}
\DeclareMathOperator*{\maxo}{max}
\DeclareMathOperator*{\erf}{erf}
\DeclareMathOperator*{\arcosh}{arcosh}
\DeclareMathOperator*{\sgn}{sign}
\DeclareMathOperator*{\forallo}{\Large\forall}
\newcommand{\ie}{{\em i.e.}\ }
\newcommand{\eg}{{\em e.g.}\ }
\newtheorem{defn}{Definition}
\newtheorem{lemma}{Lemma}
\newtheorem{cor}{Corollary}
\newtheorem{prop}{Proposition}
\newtheorem{theorem}{Theorem}
\newenvironment{proof}[1][]{\paragraph{Proof #1:}}{\hfill$\square$}
\patchcmd{\thebibliography}{\section*{\refname}}{\section*{\refname}}{}{}
\let\oabsname\abstractname
\renewcommand{\abstractname}{\large{\oabsname}}

\newenvironment{Itemize}{\vspace{-0.5\parskip}\begin{itemize}}{\end{itemize}}
\newenvironment{Enumerate}{\vspace{-0.5\parskip}\begin{enumerate}}{\end{enumerate}}
%\setlist*[itemize]{noitemsep, topsep=0pt}
\setlist*[itemize]{itemsep=5pt, topsep=0pt}
\setlist*[enumerate]{itemsep=5pt, topsep=0pt}

%\definecolor{bggrey}{rgb}{0.95,0.95,0.95}
\definecolor{bggrey}{rgb}{0.93,0.95,0.94}
%\definecolor{bggrey}{rgb}{239,241,242}

\lstdefinestyle{mystyle} {
  framexleftmargin=2pt, %keeps numbers outside of frame
  %framextopmargin=15pt,
  %framexbottommargin=-15pt,
  frame=tb,
  framesep=7pt,
  rulecolor=\color{bggrey},
  backgroundcolor=\color{bggrey},
  basicstyle=\ttfamily\scriptsize,
  linewidth=1.0\textwidth,
  numbers=left,
  showlines=true,
  numberblanklines=false,
  stepnumber=1,
  numbersep=5pt,
  %numberstyle=\ttfamily\scriptsize,
  breaklines=true,
  breakatwhitespace=true,
  tabsize=2,
  keepspaces=true,
  mathescape=true,
  escapeinside={(*@}{@*)}
}

\lstdefinestyle{mystyle1} {
  %backgroundcolor=\color{bggrey},
  %basicstyle=\ttfamily\scriptsize,
  linewidth=0.5\textwidth,
  numbers=none,
  belowskip=-0.8 \baselineskip,
  %breaklines=true,
  %breakatwhitespace=true
}
\lstset{style=mystyle}

% don't need the following. simply use defaults
\setlength{\baselineskip}{16.0pt}    % 16 pt usual spacing between lines
% Default margins are too wide all the way around.  I reset them here
\setlength{\parskip}{2ex plus 2pt}
\setlength{\parindent}{0pt}
\setlength{\oddsidemargin}{.125in}
%\setlength{\oddsidemargin}{0.5cm}
\setlength{\evensidemargin}{0.5cm}
\setlength{\marginparsep}{0.75cm}
\setlength{\marginparwidth}{2.5cm}
\setlength{\marginparpush}{1.0cm}
\setlength{\textwidth}{6.25in}
%\setlength{\textwidth}{150mm}
\setlength{\textheight}{9in}
\setlength{\topmargin}{-.5in}

\renewcommand{\headrulewidth}{0pt}
\renewcommand{\footrulewidth}{0pt}
\fancypagestyle{fpstyle}{ %define a style for fitst page
  \fancyhf{}
  \fancyfoot[L]{{\small\em Working Paper}}
  %\fancyfoot[R]{{\small\em Confidential}}
}
\fancyhf{} %clear all header and footer fields
\fancyfoot[L]{{\small\em Working Paper}}
%\fancyfoot[R]{{\small\em Confidential}}
\fancyfoot[C]{\thepage}
                 
\title{
  \textsc{Interest Rate Volatility Surface Simulation\\
  for Risk Management}\\%\footnotemark
  \large{\em{(Working Paper)}}
}
\author{Maged Tawfik%\footnotemark %\\
%Citibank
}
\renewcommand{\today}{\vspace{-20pt}\it \small{May 28, 2026}}

\begin{comment}
\end{comment}

\begin{document}
\pagestyle{fancy}

%\leftdisplays
\maketitle
\thispagestyle{fpstyle} % use if do not want page numbers
%\pagestyle{empty} % all pages are not numbered
\begin{center}
\end{center}
\footnotetext[1]{The opinions expressed in the paper are those of the author and do not necessarily reflect the view of Citibank.}
\footnotetext[2]{email: maged.tawfik@citi.com}
%\begin{abstract}
%\input{liqNetflowAbstract_txt.tex}
%\end{abstract}
\section{Problem Setup}
The implied volatility of an option on an interest rate contract
is the model volatility that makes the model price match the
market price of the option.
Swaps, forwards and future contracts
are examples of the underlying instruments.
Swaptions, caps and floors are examples of the options.
The quoted implied volatility is usually either the normal (Bachelier)
volatility, the Black (lognormal) volatility or the CEV (Constant
Elasticity of Variance) volatility.
The choice depends on the model used to price the option.
Converting from one form of implied volatility to another can be
typically done using either closed from or semi-analytic approximations.

The price function of a simple option $C$ can be assumed a function of
the form $C(v\,; F, k, x): v\mapsto c$, where $c$ is the option price,
and $F, v, k, x$ are the risk neutral forward price of the
underlying contract, its implied volatility,
the strike price of the option and
its time to expiry, respectively.
For at-the-money options, $k=F$.
Let $\mathbb{C}_v$ be the space of functions $C_v:(F,k,x)\mapsto c$,
\ie $C_v(.)=C(v;.)$, which are monotonic in $v$.
The implied volatility surface, of such options, can be formally defined as the map
\beqn
\begin{aligned}
    V\!:&\;X\times T\times K\rightarrow \mathbb{C}_v\\
    &(x,t,k)\mapsto C_v(.)\\
    &(x,t,k)\mapsto v
\end{aligned}
\eeqn
where $X$ is a set of times to expiry, $T$ is a set of contract tenors and
$K$ is a set of strike prices, and where $v$ is the implied volatility. 
Put differently, $V(x,t,k)=C^{-1}(c;F,k,x)$.

It is not uncommon, in options modeling,
to use appropriate models to relate the price of
off-the-money options to at-the-money ones of the same
expiry date.
Most such models are rich enough to handle smile and skew effects.
As a result, the dimensionality of the volatility surface can be reduced
by limiting it to at-the-money options.
In such case,
the implied volatility surface reduces to,
$V: X\times T\rightarrow \mathbb{C}_v$.

\section{Forward Stochastic model}
The current common options market practice is to model the stochastic behavior of interest rate contracts
using the SABR model, namely

\beqn\label{eq:sabr}
\begin{aligned}
&dF =\; F^\beta \alpha\, dW_F\\
&d\alpha =\;\alpha\, \nu\, dW_\alpha\\
&<dW_F,\, dW_\alpha> \;=\rho
\end{aligned}
\eeqn
where $F$ is the risk neutral forward price of the underlying contract,
and $dW_F, dW_\alpha$ are the Wiener increments associated with
the forward price and its volatility, respectively.
Notice that $\alpha$ represents the volatility of the forward contract,
while $\nu$ represents the volatility of $\alpha$, and thus
termed the vol of vol.

Semi-analytic methods, such as those given by Hagan and by Palout
(see Choi and Wu \cite{CW} for a full discussion on this topic)
are typically used to convert a SABR volatility tuple $(\alpha, \beta, \rho, \nu)$
to an equivalent normal, Black or CEV volatility.
If we assume the existence of a function
$H:(\alpha,\beta,\rho,\nu;\,x\in X, t\in T, k\in K)\mapsto v$,
then simulating the values of $\alpha$ over $X\times T$
should suffice to fully describe the dynamics of the implied
volatility surface.

The inverse of the monotonic $H(\alpha|.)$, which we shall denote $H^{-1}(v|.)$,
can be used to calculate SABR $\alpha$ from implied volatility,
given the reset of the parameters.
For a given $x$ and $t$ observable option prices, at a given point in time,
for a large enough set of $k$ values can be used
to solve for the parameters $\alpha, \beta, \rho$ and $\nu$.
The parameters $\beta, \rho$ and $\nu$ tend to remain stable for extended periods of time,
and once determined, can be assumed to remain constant during simulation.

For the current work, we will use a CIR variation of
the SABR model \eqref{eq:sabr}.
Our model, has the following dynamics
\beqn\label{eq:cir}
\begin{aligned}
&dF_{x,t} =\; F_{x,t}^{\beta_{x,t}} \alpha_{x,t}\, dW_{F_{x,t}}\\
&d\alpha_{x,t} =\;\theta_{x,t}(\mu_{x,t} - \alpha_{x,t})\,dt\,+\, \nu_{x,t}\,\sqrt\alpha_{x,t}\, dW_{\alpha_{x,t}}\\
&<\!dW_{F_{x,t}},\, dW_{\alpha_{x,t}}\!> \;=\rho_{x,t}
\end{aligned}
\eeqn
where the index $x,t$ refers to the point in the grid,
and $\mu_{x,t}$ and $\theta_{x,t}$ are the process mean and
the mean reversion speed, respectively.
Unless needed for clarity, we will drop the $x,t$ index in the remainder of the work for simplicity of presentation.

\section{Parameter Estimation}
We rely on historical time series analysis for 
the estimation of the parameters associated with the model given by \eqref{eq:cir}.
First, we construct a time series of sufficient length for the implied at-the-money volatility
of a given tenor and time to expiry.
We then use a semi-analytic model such as the modified Hagan model \cite{HKLW} or the Paulot \cite{P} model
to convert
the implied volatility time series to a time series for the SABR $\alpha$ of \eqref{eq:cir}.

Prykhodko and Ralchenko \cite{PR} suggested the following as an unbiased estimator for the diffusion parameter,
assuming ergodic behavior and given a time series of forward prices $F_i, i\le n$,
that covers a period of $t_n$ years,

\beqn
\hat{\nu}^2 = \frac{n}{t_n}\,\frac{\sum_{i<n} (F_i-F_{i-1})^2}{\sum_{i<n}F_i}\,.
\eeqn

Ben Alaya and Kebaier \cite{BK} showed that

\beqn
\frac{1}{t_n}\,\int_0^{t_n} F_s ds\,\xrightarrow{\;\;\mathbb{P}\;\;}\,\mu\quad \text{as} \quad t\!\longrightarrow\!\infty\,.
\eeqn

They also suggested the following for the consistent estimation of $\theta$

\beqn
\hat{\theta}\,=\,\frac{t_n\,(\log{F_n}-\log{F_0}+\frac{t_n}{2n}S^2\sum_{i<n}\frac{1}{F_i})-\frac{t_n}{n}(F_n-F_0)\sum_{i<n}\frac{1}{F_i}}
{t_n^2\,(\,\frac{1}{n^2}\sum_{i<n}\frac{1}{F_i}\,\sum_{i<n}F_i\;-\;1\,)}\,.
\eeqn

\section{PCA Analysis}
Detailed time series analysis of SOFR interest rate swaptions between 2020 and 2026, showed significant correlation.
The options included in the dataset included 17 times to expiry and 23 tenors, \ie
the $X\times T$ set had a cardinality of 391.
It was found out that 76\% of the total variance in the
normalized (centered and scaled) daily change in the SABR alpha value, was explained by the first principal component.
The first three components were sufficient to explain 94\% of the total variance.
Whereas, 98\% was explained by the first five components.
These results suggest that the full implied volatility surface can
be reliably simulated using a handful of stochastic factors.

Given a historical time series matrix $M$ of $N_h$ rows and
$X\times T$ columns, where $N_h$ is the number of observations of $\alpha$ values for each grid point,
the SVD (Singular Value Decomposition) of X is given by
\beqn
M = U\;S\;V^\top
\eeqn
where $M$ is an $N_h\times(X\times T)$ matrix, $U$ an $N_h\times N_h$ matrix,
$S$ an $N_h\times(N\times T)$ diagonal matrix of non-negative non-increasing real numbers,
and $V$ an $(X\times T)\times(X\times T)$ matrix.

Retaining only the top $p$ principal components, the SABR $\alpha_{x,t}$ of any option $(x,t)$ can be expressed
as

\beqn
\begin{aligned}
d\alpha_{x,t}=&\theta_{x,t}(\mu_{x,t} - \alpha_{x,t})\,dt\,+\, \nu_{x,t}\,\sqrt\alpha_{x,t}\,
\sum_{i\le p}\rho_{i,x,t}\, dW_i\\
\rho_{i,x,t}=&\frac{1}{R}\,S_{i,i}\,V_{(x,t),i}\\[10pt]
R=&\sqrt{\sum_{i\le p}(S_{i,i}V_{(x,t),i})^2}
\end{aligned}
\eeqn
where $dW_i$ are $p$ Wiener increments for each time step.

\section{Monte Carlo Simulation}
The nonlinearity involved in the CIR process requires delicate attention to the discretization
scheme employed in the Monte Carlo simulation.
The algorithms developed by Alfonsi \cite{A} offer a few higher order reliable schemes,
which allow a larger simulation time step without loss of accuracy while preserving the
non-negativity condition on $\alpha$.

In the current work, we chose the $E(0)$ discretization algorithm specified by Alfonsi \cite{A}.
Using our notation, Alfonsi's $E(0)$ is given by

\beqn
\hat{\alpha}_{t_{i+1}}=\hat{\alpha}_{t_i}(1-\theta \delta)\,+\,\nu\sqrt{\hat{\alpha_{t_i}}}(W_{t_{i+1}}-W_{t_i})\,+
\frac{1}{4}\nu^2(W_{t_{i+1}}-W_{t_i})^2 + \delta(\mu\theta-\frac{1}{4}\nu^2)
\eeqn
where $\delta$ is the simulation time step.

\begin{thebibliography}{99}
    \bibitem{A}
        Aurélien, Alfonsi, 2010, "High Order Discretization Schemes
        For The Cir Process: Application
        To Affine Term Structure And Heston Models,"
        Mathematics Of Computation, 79(269), 209-237
    \bibitem{BK}
        Ben Alaya M, Kebaier A, 2013, “Asymptotic Behavior of the Maximum Likelihood Estimator for Ergodic and Nonergodic Square-root Diffusions.” Stochastic Analysis and Applications, 31(4), 552-573. doi:10.1080/07362994.2013.798175.
    \bibitem{CW}
        Choi, Jaehyuk and Wu, Lixin, 2020, "The Equivalent Constant-elasticity-of-variance (CEV) Volatility
        of the Stochastic-Alpha-Beta-Rho (SABR) Model,"
        arXiv:1911.13123v3
    \bibitem{HKLW}
        Hagan, Patrick; Kumar, Deep; Lesnieweski, Andrew S. and Woodward, Diana, 2014,
        "Arbitrage-Free SABR," {\em Wilmott} (2014):60-75. doi:10.1002/wilm.10290
    \bibitem{P}
        Paulot, Louis, 2015, "Asymptotic Implied Volatility at the Second Order with
        Application to the SABR Model," In: Friz, P., Gatheral, J., Gulisashvili,
        A., Jacquier, A., Teichmann, J. (eds) {\em
        Large Deviations and Asymptotic Methods in Finance}.
        Springer Proceedings in Mathematics \& Statistics, (110)
        Springer, Cham. https://doi.org/10.1007/978-3-319-11605-1\_2
        arXiv:0906.0658v3
    \bibitem{PR}
        Prykhodko, Olha, and Ralchenko, Kostiantyn, 2025, "Discretization and Asymptotic Normality of Drift Parameters Estimator
        in the Cox-Ingersoll-Ross Model," Austrian Journal of Statistics, 54, 82-99. doi:10.17713/ajs.v54i1.1968
\end{thebibliography}
\end{document}